\documentclass[12pt]{article}

\input{../../_assets/preambulo.tex}

\begin{document}

    % 1. Foto de fondo
    % 2. Título
    % 3. Encabezado Izquierdo
    % 4. Color de fondo
    % 5. Coord x del titulo
    % 6. Coord y del titulo
    % 7. Fecha

    \input{../../_assets/portada}
    \portadaExamen{ffccA4.jpg}{Variable Compleja I\\Examen XIX}{Variable Compleja I. Examen XIX}{MidnightBlue}{-9.5}{28}{2025-26}{David Muñoz Gómez}

    \begin{description}
        \item[Asignatura] Variable Compleja I.
        \item[Curso Académico] 2025/2026.
        \item[Grado] Ingeniería Informática y Matemáticas.
        \item[Grupo] Único.
        \item[Profesor] Javier Meri de la Maza.
        \item[Descripción] Convocatoria ordinaria.
        \item[Fecha] 19 de junio de 2026.
        \item[Duración] 3h.
    \end{description}
    \newpage

    % ------------------------------------
    
    \begin{ejercicio}[1.5 puntos]
        Sea $R \in \mathbb{R}^+$ y $f \in \mathcal{H}(D(0, R))$, no constante. Pruebe que la función $M \colon \mathopen{]}0, R\mathclose{[} \longrightarrow \mathbb{R}$ definida por
        \[ M(r) = \max \{|f(z)| : z \in C(0, r)^*\} \quad \forall r \in \mathopen{]}0, R\mathclose{[} \]
        es estrictamente creciente.
    \end{ejercicio}

    \begin{ejercicio}[2 puntos]
        Sea $f$ una función entera verificando que existen constantes $R, \veps > 0$ tales que
        \[ |f(z)| \ge \veps, \text{ si } |z| \ge R. \]
        Pruebe que $f$ es constante o tiene algún cero en $D(0, R)$.
    \end{ejercicio}

    \begin{ejercicio}[2.5 puntos]
        Integrando una conveniente función compleja sobre la poligonal $[-R, R, R + \pi i, -R + \pi i, -R]$ con $R \in \mathbb{R}^+$, calcule la integral
        \[ \int_{-\infty}^{\infty} \frac{\cos(2x)}{e^x + e^{-x}} \, dx. \]
    \end{ejercicio}

    \begin{ejercicio}[1 punto por cada apartado]
        Demuestre los siguientes enunciados:
        \begin{enumerate}
            \item Sea $f \in \mathcal{H}(\mathbb{C})$ verificando $f\left(\nicefrac{1}{n}\right) = \nicefrac{1}{n^3}$ para cada $n \in \mathbb{N}$. Calcule $f$.
            \item La siguiente función es holomorfa en $\bb{C}\setminus \{0\}$ y admite una extensión entera.
            \Func{f}{\bb{C}\setminus \{0\}}{\bb{C}}{z}{\displaystyle \int_{C(0,1)} \frac{e^{zw} - 1}{zw} \, dw}
            \item No existe una función $f \in \mathcal{H}(\mathbb{C})$ tal que $|z^2| + |e^z| \le |f(z)|$ para cada $z \in \mathbb{C}$.
            \item Dé un ejemplo de función $f \in \mathcal{H}(D(0, 1))$ no acotada con $\text{Im}(f)$ acotada.
        \end{enumerate}
    \end{ejercicio}

    \newpage
    \noindent
    \textbf{Solución.}
    \setcounter{ejercicio}{0} % Reiniciar contador de ejercicios

    \begin{ejercicio}[1.5 puntos]
        Sea $R \in \mathbb{R}^+$ y $f \in \mathcal{H}(D(0, R))$, no constante. Pruebe que la función $M \colon \mathopen{]}0, R\mathclose{[} \longrightarrow \mathbb{R}$ definida por
        \[ M(r) = \max \{|f(z)| : z \in C(0, r)^*\} \quad \forall r \in \mathopen{]}0, R\mathclose{[} \]
        es estrictamente creciente.
        
        \bigskip
        Dados $r_1, r_2 \in \mathopen{]}0,R\mathclose{[}$ con $r_1 < r_2$, veamos que $M(r_1) < M(r_2)$. Como $r_2 < R$, $\overline{D}(0,r_2) \subset D(0,R)$. Por el Principio del Módulo Máximo, tenemos que:
        \[ \max\{|f(z)| : z \in \overline{D}(0,r_2)\} = \max\{|f(z)| : z \in C(0,r_2)^*\} = M(r_2). \]
        Además, la desigualdad es estricta para todo punto del interior, es decir, si $z \in~D(0,r_2)$ entonces $|f(z)| < M(r_2)$, pues el Principio del Módulo Máximo indica que en caso de darse la igualdad en algún punto interior, $f$ es constante. Pero $f$ no lo es por hipótesis.

        Como $r_1 < r_2$, tenemos que $\overline{D}(0,r_1) \subset D(0,r_2)$. En particular, si $z \in C(0,r_1)^* \subset D(0,r_2)$, se cumple que $|f(z)| < M(r_2)$.\\

        Por lo tanto, tomando el máximo sobre la circunferencia de radio $r_1$, resulta:
        \[ \max\{|f(z)| : z \in C(0,r_1)^*\} = M(r_1) < M(r_2) \]
    \end{ejercicio}

    \begin{ejercicio}[2 puntos]
        Sea $f$ una función entera verificando que existen constantes $R, \veps > 0$ tales que
        \[ |f(z)| \ge \veps, \text{ si } |z| \ge R. \]
        Pruebe que $f$ es constante o tiene algún cero en $D(0, R)$.
        
        \bigskip

        Supongamos que $f$ no es una función polinómica. En ese caso, por el corolario del Teorema de Casorati sabemos que:
        \[ \exists \{z_n\} \rightarrow \infty\qquad \text{tal que} \qquad \{f(z_n)\} \rightarrow 0 \]

        Por tanto, por la divergencia de la sucesión $\{z_n\}$ está garantizado que 
        \[\exists n_1 \in \mathbb{N} :\; \forall m \geq n_1 \Rightarrow |z_m| \geq R\]

        Por la convergencia de la sucesión $\{f(z_n)\}$ está garantizado que 
        \[\exists n_2 \in \mathbb{N} : \; \forall m \geq n_2 \Rightarrow |f(z_m)| < \veps\]

        Tomando $n=\max\{n_1,n_2\}$ vemos que $\forall m \geq n : \; |f(z_m)| < \veps \wedge |z_m| \geq R$\\
        
        Esto es una contradicción con la propiedad que verifica $f$.\\

        Por tanto, $f$ debe ser un polinomio o bien constante o de grado mayor o igual que 1. Supongamos que $f$ no es constante, entonces como hemos dicho $\deg(f) \geq 1$.\\

        Por el Teorema fundamental del Álgebra: $\exists z\in \mathbb{C} : \; f(z) = 0$, aunque realmente existirán tantos como el grado de f contando multiplicidad. Como $f(z)=0 \Rightarrow |f(z)| = 0 < \veps$, tenemos que $|z| < R \iff z\in D(0,R)$\\

        Por tanto, concluimos que o $f$ es un polinomio constante o tiene una raíz en $D(0,R)$.
    \end{ejercicio}

    \begin{ejercicio}[2.5 puntos]
        Integrando una conveniente función compleja sobre la poligonal $[-R, R, R + \pi i, -R + \pi i, -R]$ con $R \in \mathbb{R}^+$, calcule la integral
        \[ \int_{-\infty}^{\infty} \frac{\cos(2x)}{e^x + e^{-x}} \, dx. \]
        
        % AQUÍ IRÁ LA SOLUCIÓN
    \end{ejercicio}

    \begin{ejercicio}[1 punto por cada apartado]
        Demuestre los siguientes enunciados:
        \begin{enumerate}
            \item Sea $f \in \mathcal{H}(\mathbb{C})$ verificando $f\left(\nicefrac{1}{n}\right) = \nicefrac{1}{n^3}$ para cada $n \in \mathbb{N}$. Calcule $f$.
            
                \bigskip

                Sea $A = \{\nicefrac{1}{n}:\;n \in \mathbb{N}\}$. Claramente $A$ tiene al $0$ como punto de acumulación.

                Como $f \in \mathcal{H}(\mathbb{C})$ y $A'\cap \mathbb{C} \neq \emptyset$, aplicando el principio de identidad:
                \[f(z) = z^3 \ \forall z\in A \qquad \Longrightarrow \qquad f(z) = z^3 \ \forall z \in \mathbb{C}\]

            \item La siguiente función es holomorfa en $\bb{C}\setminus \{0\}$ y admite una extensión entera.
            \Func{f}{\bb{C}\setminus \{0\}}{\bb{C}}{z}{\displaystyle \int_{C(0,1)} \frac{e^{zw} - 1}{zw} \, dw}

                \bigskip
                Sea la siguiente función:
                \Func{\phi}{C(0,1)\times \bb{C}\setminus \{0\}}{\bb{C}}{(w,z)}{\dfrac{e^{zw}-1}{zw}}

                Como $zw \neq 0$ para todo $w\in C(0,1), z \in \mathbb{C}\setminus\{0\}$ vemos que $\phi(w,z)$ es continua.
                
                Además, $\phi_w(z) = \phi(w,z)$ es cociente de funciones holomorfas en $\mathbb{C}\setminus\{0\}$ luego $\phi_w(z) \in \mathcal{H}(\mathbb{C}\setminus\{0\})$. Por el Teorema de integrales dependientes de parámetros $f \in \mathcal{H}(\mathbb{C}\setminus\{0\})$.

                \bigskip

                Para ver que tiene extensión entera hay muchas formas de hacerlo, pero una manera curiosa y con resultados inesperados es como sigue. Calculamos explícitamente el valor de la integral para $z \in \mathbb{C} \setminus \{0\}$. 
                Separando el integrando por linealidad vemos que:
                \[f(z) = \frac{1}{z} \int_{C(0,1)} \frac{e^{zw}}{w} \, dw - \frac{1}{z} \int_{C(0,1)} \frac{1}{w} \, dw\]

                Aplicamos la Fórmula Integral de Cauchy en $w_0 = 0$ a la función entera $g(w) = e^{zw}$:
                \[\int_{C(0,1)} \frac{e^{zw}}{w} \, dw = 2\pi i \cdot g(0) = 2\pi i\]

                Sabiendo por parametrización usual que $\displaystyle \int_{C(0,1)} \frac{1}{w} \, dw = 2\pi i$ obtenemos:
                \[f(z) = \frac{1}{z} (2\pi i - 2\pi i) = 0 \quad \forall z \in \mathbb{C}\setminus\{0\}\]

                Para ver el comportamiento en el origen calculamos el siguiente límite:
                \[\lim_{z \to 0} f(z) = \lim_{z \to 0} 0 = 0\]

                Como el límite existe y es un valor finito en $\mathbb{C}$, por el Teorema de Extensión de Riemann la singularidad en $z=0$ es evitable.
                Por tanto, $f$ admite una extensión entera $\tilde{f} \in \mathcal{H}(\mathbb{C})$ dada por $\tilde{f}(z) = 0 \; \forall z \in \mathbb{C}$.


            \item No existe una función $f \in \mathcal{H}(\mathbb{C})$ tal que $|z^2| + |e^z| \le |f(z)|$ para cada $z \in \mathbb{C}$.
                
                \bigskip

                Supongamos que f es no polinomica. Como además es entera, por el corolario del Teorema de Casorati:
                \[ \exists \{z_n\} \rightarrow \infty\qquad \text{tal que} \qquad \{f(z_n)\} \rightarrow 0 \]

                Pero como $z^2$ es un polinomio, $|z_n^2| \rightarrow \infty$ y $|e^z_n| > 0$ pero $|z_n^2| + |e^z_n| \le |f(z_n)|$.

                Esto es una contradicción, por tanto $f$ solo puede ser un polinomio, pero si $f$ es un polinomio:

                \[|z^2| + |e^z| \leq |f(z)| \Rightarrow |e^z| \leq |f(z)| \; \forall z\in\mathbb{C}\]

                Esta desigualdad nos indica que la función $e^z$ tiene crecimiento subpolinómico no solo en el exterior de un disco\footnote{Ejercicio 2 de la relación de ceros de funciones holomorfas.} sino en todo $\mathbb{C}$, en cuyo caso $e^z$ sería un polinomio.
            \item Dé un ejemplo de función $f \in \mathcal{H}(D(0, 1))$ no acotada con $\text{Im}(f)$ acotada.
            
                \bigskip

                Sea la siguiente función, que emplea el logaritmo principal:
                \Func{f}{\mathbb{C}\setminus\mathopen{]}-\infty,-1\mathclose{]}}{\mathbb{C}}{z}{\log(1+z)}

                Como $\log(1+z) = \ln|1+z| + i\arg(1+z)$ y el $\ln$ es una función estrictamente creciente y $|\arg(1+z)| \leq \pi$ este es un ejemplo de función con imagen no acotada pero parte imaginaria acotada.
        \end{enumerate}
    \end{ejercicio}

\end{document}